\section{Introduction}

Following the outbreak of a public emergency, social media quickly becomes a central space for both information dissemination and emotional expression. Public opinion often evolves through successive stages, including outbreak, clarification, reactivation, and decline, and its trajectory is jointly shaped by collective psychology, official information release, and platform mechanisms. Understanding these complex dynamics is therefore essential for crisis management and decision-making.

Early computational approaches mainly rely on mathematical models to simulate the dynamics of public opinion. Classical methods, such as epidemic models (SIR/SIS), the DeGroot consensus model, and the Deffuant bounded-confidence model, can capture macroscopic patterns, but they reduce individual expressions to simple numerical states and therefore cannot represent the complexity of natural language interaction or emotional dynamics. With the rapid progress of large language models (LLMs), high-fidelity social simulation becomes increasingly feasible. LLM-driven agents can generate semantically rich behavioral outputs and show clear advantages in role playing and contextual understanding.

Current studies on multi-agent public opinion simulation exhibit a pronounced risk-diffusion orientation in their problem settings. Their main focus is on scenarios such as rumor propagation, misinformation diffusion, emotional polarization, information-pollution attacks, and opinion manipulation. For example, RumorSphere concentrates on the dynamic process of rumor spread~\citep{liu2025rumorsphere}, whereas TrendSim examines topic evolution under trend manipulation and adversarial attacks~\citep{zhang2025trendsim}. Although these studies advance the scale of simulation and the quality of generated behaviors, their research frameworks are still largely centered on abnormal diffusion or adversarial settings. In addition, most existing frameworks advance simulation with fixed rounds, which makes it difficult to capture the asynchronous participation of real social media users and the stage-wise evolution of real-world events. Many studies also emphasize end-to-end output quality while lacking ablation designs that compare individual modules against real event data, making it difficult to isolate the independent contribution of each design component to public opinion dynamics. More importantly, the vast majority of prior work is validated on only a single case, and thus lacks cross-event generalization tests, which limits the external validity of the resulting conclusions.

To address these limitations, this paper develops a time-aware multi-agent public opinion simulation system for public emergencies and validates it on three real emergency events of different types. The main contributions are threefold. First, we propose a simulation framework that integrates continuous-time scheduling, official publisher injection, and a cognitive reflection mechanism. Second, we verify the independent contributions of the official publisher and reflection mechanism through systematic ablation experiments. Third, we analyze the heterogeneous emotional response effects induced by official publisher events across different cases.
